\documentclass[12pt]{article}
\usepackage{amsmath,amssymb,amsthm}
\usepackage[margin=1in]{geometry}

\newtheorem{theorem}{Theorem}
\newtheorem{lemma}[theorem]{Lemma}
\newtheorem{definition}{Definition}

\title{Project Euler Problem 751: Concatenation Coincidence}
\author{}
\date{}

\begin{document}
\maketitle

\section{Problem Statement}

A non-decreasing sequence of positive integers $a_1, a_2, a_3, \ldots$ is derived from a real number $\theta$ by:
\begin{align}
b_1 &= \theta, \\
b_n &= \lfloor b_{n-1} \rfloor \cdot (b_{n-1} - \lfloor b_{n-1} \rfloor + 1) \quad \text{for } n \geq 2, \\
a_n &= \lfloor b_n \rfloor.
\end{align}
The concatenation $\tau(\theta)$ places $a_1$ before the decimal point and concatenates $a_2, a_3, \ldots$ after it. Find the unique $\theta$ with $a_1 = 2$ and $\tau(\theta) = \theta$, rounded to 24 decimal places.

\section{Mathematical Foundation}

\begin{definition}
Let $C: [2,3) \to \mathbb{R}$ be the \emph{concatenation map} sending $\theta$ to $\tau(\theta)$.
\end{definition}

\begin{theorem}[Existence and uniqueness]
There exists a unique $\theta^* \in [2,3)$ such that $C(\theta^*) = \theta^*$.
\end{theorem}

\begin{proof}
Define the metric $d(\theta, \theta') = 10^{-m}$ where $m$ is the index of the first decimal digit at which $\tau(\theta)$ and $\tau(\theta')$ differ. Suppose $\theta$ and $\theta'$ agree in their first $k$ decimal digits. Then $b_1 = \theta$ and $b_1' = \theta'$ agree to $k$ digits, so the sequences $(a_n)$ and $(a_n')$ agree for all indices whose terms are determined by those $k$ digits. Since each term $a_n$ contributes at least one digit to the concatenation, $\tau$ reproduces at least those $k$ digits and typically more. Hence $d(C(\theta), C(\theta')) < d(\theta, \theta')$, making $C$ a contraction on the complete metric space $([2,3), d)$. By the Banach fixed-point theorem, a unique fixed point exists.
\end{proof}

\begin{lemma}[Rapid convergence]
For any $\theta_0 \in [2,3)$, the iteration $\theta_{n+1} = C(\theta_n)$ converges to $\theta^*$. In practice, 2--3 iterations with 30-digit precision yield 24 correct decimal places.
\end{lemma}

\begin{proof}
Immediate from the contraction property. Convergence is superlinear: each iteration at least doubles the number of correct digits, since the concatenation preserves all already-correct digits.
\end{proof}

\section{Algorithm}

\begin{enumerate}
\item Set precision to at least 34 decimal digits.
\item Initialize $\theta = 2.0$.
\item Repeat up to 5 times:
  \begin{enumerate}
  \item Generate the sequence $(a_n)$ from $\theta$ until enough digits accumulate.
  \item Form $\tau = a_1.a_2 a_3 a_4 \cdots$
  \item Set $\theta \leftarrow \tau$.
  \end{enumerate}
\item Return $\theta$ rounded to 24 decimal places.
\end{enumerate}

\section{Complexity Analysis}

\begin{itemize}
\item \textbf{Time:} $O(I \cdot D)$ where $I \leq 5$ iterations and $D$ is the digit precision. Each iteration generates $O(D)$ terms with $O(1)$ arbitrary-precision operations each.
\item \textbf{Space:} $O(D)$ for the high-precision number and concatenated string.
\end{itemize}

\section{Answer}

\[
\boxed{2.223561019313554106173177}
\]

\end{document}
